% ============================================================
%  General Mathematical Notation Table
%  Paper: One Step Diffusion via Shortcut Models (Frans et al., ICLR 2025)
%  Capstone: Statistical Machine Learning
% ============================================================

\section*{General Mathematical Notation}
\addcontentsline{toc}{section}{General Mathematical Notation}
\label{sec:notation}

The following table collects all symbols and operators used throughout this
report, organized by category.

\renewcommand{\arraystretch}{1.5}

% ── 1. Data and Distributions ────────────────────────────────────────────────
\begin{table}[h]
\centering
\caption*{\textbf{Table N1.} Data and Distributions}
\begin{tabular}{@{} c p{10cm} @{}}
\toprule
\textbf{Symbol} & \textbf{Definition} \\
\midrule
$x_0$
  & Clean (real) data sample drawn from the data distribution;
    $x_0 \sim D$. \\
$x_1$
  & Pure Gaussian noise sample;
    $x_1 \sim \mathcal{N}(0, I)$ (the ``noisy'' endpoint in flow-matching
    convention where $t=1$ is full noise). \\
$x_t$
  & Interpolated (noisy) data point at noise level $t$, defined by the
    linear interpolation:
    $x_t = (1-t)\,x_0 + t\,x_1$,\quad $t \in [0,1]$. \\
$\varepsilon$
  & Standard Gaussian noise; $\varepsilon \sim \mathcal{N}(0, I)$. \\
$D$
  & Training dataset (e.g.\ CelebA-HQ 256 or ImageNet-256). \\
$p_{\mathrm{data}}$
  & True data distribution; $x_0 \sim p_{\mathrm{data}}$. \\
$p_t$
  & Marginal distribution of $x_t$ at noise level $t$. \\
$p(d,\,t)$
  & Joint sampling distribution over step size $d$ and noise level $t$
    used during training. \\
\bottomrule
\end{tabular}
\end{table}

% ── 2. Time and Step Variables ───────────────────────────────────────────────
\begin{table}[h]
\centering
\caption*{\textbf{Table N2.} Time and Step Variables}
\begin{tabular}{@{} c p{10cm} @{}}
\toprule
\textbf{Symbol} & \textbf{Definition} \\
\midrule
$t \in [0,1]$
  & Continuous noise level (timestep).
    $t=0$: clean data; $t=1$: pure noise. \\
$d \in [0,1]$
  & Step size (\textit{shortcut conditioning parameter}).
    Controls the length of the denoising ``jump'' from $x_t$ to
    $x_{t-d}$ in a single forward pass.
    $d=0$ recovers the standard flow-matching velocity target. \\
$M$
  & Total number of denoising steps at inference
    (Number of Function Evaluations, NFE). \\
$n \in \{0,\ldots,M-1\}$
  & Inference step index. \\
$\Delta t = 1/M$
  & Uniform step increment; $d = \Delta t$ in the Euler sampling loop. \\
$d_{\mathrm{flow}}$
  & Integer encoding of step size used as conditioning input to the
    network:
    $d_{\mathrm{flow}} = \lfloor \log_2 M \rfloor \in \mathbb{Z}_{\geq 0}$.
    Reflects the binary-tree depth of the bootstrapping target
    construction. \\
\bottomrule
\end{tabular}
\end{table}

% ── 3. Model and Parameters ──────────────────────────────────────────────────
\begin{table}[h]
\centering
\caption*{\textbf{Table N3.} Model and Parameters}
\begin{tabular}{@{} c p{10cm} @{}}
\toprule
\textbf{Symbol} & \textbf{Definition} \\
\midrule
$\theta$
  & Learnable parameters of the Shortcut network. \\
$s_\theta(x,\,t,\,d)$
  & Shortcut network: maps a noisy state $x$ at noise level $t$ and
    step size $d$ to a predicted velocity (direction and magnitude of
    the denoising step). \\
$F_\theta(x_t,\,t,\,d)$
  & Equivalent shortcut mapping notation;
    $F_\theta(x_t, t, d) \approx x_{t-d}$ (predicted state after one
    shortcut step of size $d$). \\
$v_t$
  & Ground-truth velocity target in Conditional Flow Matching:
    $v_t = \frac{dx_t}{dt} = x_1 - x_0$. \\
$s_t$
  & Predicted velocity at the current state: $s_t = s_\theta(x_t, t, d)$. \\
$s_{t+d}$
  & Predicted velocity after one small step:
    $s_{t+d} = s_\theta(x_{t+d},\, t+d,\, d)$. \\
$s_{\mathrm{target}}$
  & Training regression target; set to $v_t$ (flow-matching elements)
    or $\mathrm{stopgrad}(s_t + s_{t+d})/2$ (shortcut elements). \\
$\mathrm{stopgrad}[\cdot]$
  & Stop-gradient operator: treats its argument as a constant during
    backpropagation; prevents bootstrapping collapse. \\
$\lambda$
  & Scalar loss weighting coefficient balancing the flow-matching
    and shortcut consistency objectives. \\
\bottomrule
\end{tabular}
\end{table}

% ── 4. Training Objectives ───────────────────────────────────────────────────
\begin{table}[h]
\centering
\caption*{\textbf{Table N4.} Training Objectives}
\begin{tabular}{@{} c p{10cm} @{}}
\toprule
\textbf{Symbol} & \textbf{Definition} \\
\midrule
$\mathcal{L}_{\mathrm{flow}}$
  & Flow-matching loss (applied to the first $k$ batch elements,
    where $d \leftarrow 0$):
    \[
      \mathcal{L}_{\mathrm{flow}}
        = \mathbb{E}_{t,\,x_0,\,\varepsilon}
          \bigl[\|s_\theta(x_t,\,t,\,0) - v_t\|^2\bigr].
    \] \\
$\mathcal{L}_{\mathrm{shortcut}}$
  & Self-consistency loss (applied to remaining batch elements):
    \[
      \mathcal{L}_{\mathrm{shortcut}}
        = \mathbb{E}_{t,\,d,\,x_0,\,\varepsilon}
          \Bigl[\bigl\|s_\theta(x_t,\,t,\,2d)
            - \mathrm{stopgrad}
              \tfrac{s_t + s_{t+d}}{2}\bigr\|^2\Bigr].
    \] \\
$\mathcal{L}$
  & Combined training loss:
    $\mathcal{L} = \mathcal{L}_{\mathrm{flow}}
                 + \lambda\,\mathcal{L}_{\mathrm{shortcut}}$. \\
$\nabla_\theta$
  & Gradient with respect to network parameters $\theta$. \\
\bottomrule
\end{tabular}
\end{table}

% ── 5. Inference (Sampling) ──────────────────────────────────────────────────
\begin{table}[h]
\centering
\caption*{\textbf{Table N5.} Inference (Sampling)}
\begin{tabular}{@{} c p{10cm} @{}}
\toprule
\textbf{Symbol} & \textbf{Definition} \\
\midrule
$x \sim \mathcal{N}(0,I)$
  & Initial noise sample at $t=0$ (start of the reverse process). \\
$x \leftarrow x + s_\theta(x,\,t,\,d)\cdot d$
  & Euler update rule: advances the state by one step of size $d$
    along the predicted velocity field (Algorithm 2 of the paper). \\
$\mathrm{NFE}$
  & Number of Function Evaluations: total number of network forward
    passes used during one sampling trajectory
    ($\mathrm{NFE} = M$). \\
$\mathrm{EMA}$
  & Exponential Moving Average of network parameters used to
    stabilize training; inference always uses EMA weights. \\
\bottomrule
\end{tabular}
\end{table}

% ── 6. Evaluation Metric (FID) ───────────────────────────────────────────────
\begin{table}[h]
\centering
\caption*{\textbf{Table N6.} Evaluation Metric — Fréchet Inception Distance}
\begin{tabular}{@{} c p{10cm} @{}}
\toprule
\textbf{Symbol} & \textbf{Definition} \\
\midrule
$\mathrm{FID}$
  & Fréchet Inception Distance: scalar measure of distributional
    similarity between generated and reference images
    (lower is better). \\
$D_{\mathrm{inc}} = 2048$
  & Dimensionality of the Inception-v3 feature space used for FID
    computation. \\
$\mu_r,\,\Sigma_r$
  & Mean vector and covariance matrix of Inception features extracted
    from the \emph{reference} image set
    ($\mu_r \in \mathbb{R}^{2048}$,\
     $\Sigma_r \in \mathbb{R}^{2048 \times 2048}$). \\
$\mu_g,\,\Sigma_g$
  & Mean vector and covariance matrix of Inception features extracted
    from the \emph{generated} image set. \\
$\mathrm{FID}(\cdot\|\cdot)$
  & FID formula:
    \[
      \mathrm{FID}
        = \|\mu_r - \mu_g\|^2
          + \mathrm{tr}\!\Bigl(
              \Sigma_r + \Sigma_g
              - 2\bigl(\Sigma_r\Sigma_g\bigr)^{1/2}
            \Bigr).
    \] \\
$\mathrm{tr}(\cdot)$
  & Matrix trace operator: sum of diagonal elements. \\
$(\Sigma_r\Sigma_g)^{1/2}$
  & Matrix square root of the product $\Sigma_r\Sigma_g$;
    computed via eigendecomposition. \\
$\mathrm{FID}\text{-}50\mathrm{k}$
  & Standard absolute FID computed over 50\,000 generated samples
    against real dataset statistics (used by the original paper). \\
$\mathrm{FID}_{\mathrm{rel}}$
  & Relative FID used in this investigation: generated samples
    ($\mathrm{NFE} \in \{1,\ldots,64\}$, $n=504$) are compared
    against 128-step shortcut samples as a proxy reference,
    rather than against real data. \\
$n$
  & Number of generated samples used per FID evaluation
    ($n = 504$ in this work; $n = 50\,000$ in the original paper). \\
\bottomrule
\end{tabular}
\end{table}

\renewcommand{\arraystretch}{1}  % reset to default after notation tables
